\section{Limitations and Conclusion}

First, the historical atlas has no common structural-opportunity endpoint. Its
value is breadth and failure coverage, not causal identification across all
versions; A1 imports historical metrics and retains the historical V21
dataset-selection bias. Second, the prospective evidence is a six-dataset V21
case study with known-$K$ evaluation; its conditional signs cannot be generalized
to every structural intervention or domain. Third, E2 and E3 are diagnostics,
with dataset-by-seed aggregation and post-hoc label-aware descriptors; E3 had no
eligible replay rows, and neither stage proves objective conflict, readout
failure, or a universal local-to-global law. Fourth, the independent holdout is
zero of six and \texttt{inconclusive\_not\_completed}, so independent validation
remains unestablished; the frozen candidate/adapter manifest and its pool
shortfall were fixed before outcomes, while the recorded CUDA failure is only
incomplete compute. Finally, A2 retained E1 but prohibited E4, new Gates/losses,
and V26 or rescue routes; no new architecture or comparison was used to make the
narrative more attractive.

\paragraph{Conclusion.}
Structural information is not synonymous with clustering utility. The audited
history provides an observational atlas of heterogeneous intervention outcomes,
and a matched V21 case study shows that topology-dependent selection can be
positive or negative across data sets even when generic intervention and
selection noise are controlled. The defensible contribution is therefore not
another Gate; it is a mechanism-localization and claim-firewall protocol that
identifies where structural information may be lost and states clearly what the
current evidence does not establish.
